\hypertarget{v11-coupled_growth_full-coupled-contraction-functions-and-their-perron-growth-law}{%
\section{Coupled contraction functions and their Perron growth
law}\label{v11-coupled_growth_full-coupled-contraction-functions-and-their-perron-growth-law}}

\hypertarget{v11-coupled_growth_full-definition-and-exact-coefficient-construction}{%
\subsection{Definition and exact coefficient
construction}\label{v11-coupled_growth_full-definition-and-exact-coefficient-construction}}

Fix real parameters

\[
0<q,r,a,b<1.
\]

Consider the normalized pair

\[
\boxed{\begin{aligned}
U'(z)&=1+U(qz)V(rz),\\
V'(z)&=1+U(az)V(bz),\\
U(0)&=V(0)=0.
\end{aligned}}
\]

\textbf{Theorem 1.} There is a unique pair of entire solutions. Both
functions are odd and have strictly positive coefficients at every odd
degree. Write

\[
U(z)=\sum_{n\ge0}u_nz^{2n+1},\qquad
V(z)=\sum_{n\ge0}v_nz^{2n+1}.
\]

Then \(u_0=v_0=1\), and, for \(n\ge1\),

\[
\begin{aligned}
u_n&=\frac1{2n+1}\sum_{i=0}^{n-1}
 q^{2i+1}r^{2(n-1-i)+1}u_i v_{n-1-i},\\
v_n&=\frac1{2n+1}\sum_{i=0}^{n-1}
 a^{2i+1}b^{2(n-1-i)+1}u_i v_{n-1-i}.
\end{aligned}
\]

\textbf{Proof.} Comparing successive coefficients in the differential
equations determines the coefficient of degree \(k+1\) from coefficients
of degrees at most \(k\). Thus the formal pair is unique. The displayed
recurrence constructs an odd formal pair satisfying both equations.
Formal uniqueness excludes any additional even coefficients. Strict
positivity follows by induction.

The elementary bounds

\[
0<u_n\le2^{-n},\qquad 0<v_n\le2^{-n}
\]

hold by simultaneous induction. If they hold below \(n\), each product
in either recurrence is at most \(2^{-(n-1)}\), and there are \(n\)
products. The resulting bound is \(n2^{-(n-1)}/(2n+1)<2^{-n}\). Hence
the two series converge on \(|z|<\sqrt2\), and define a holomorphic
solution there.

Put \(Q=\max(q,r,a,b)<1\). If the pair is already defined on \(|z|<R\),
then

\[
\int_0^z[1+U(qw)V(rw)]\,dw,\qquad
\int_0^z[1+U(aw)V(bw)]\,dw
\]

define holomorphic extensions on \(|z|<R/Q\). The integrands are
holomorphic on that disk. On the old disk the new functions agree with
the old pair because their derivatives and values at zero agree. On the
larger disk, every contracted argument still lies in the old disk, so
the extensions satisfy the same equations. Iterating exhausts the
complex plane. Any other entire solution agrees with the unique germ and
therefore everywhere. \(\square\)

In particular,

\[
\begin{aligned}
U(z)&=z+\frac{qr}{3}z^3
 +\frac{qr(qr\,q^2+ab\,r^2)}{15}z^5+O(z^7),\\
V(z)&=z+\frac{ab}{3}z^3
 +\frac{ab(qr\,a^2+ab\,b^2)}{15}z^5+O(z^7).
\end{aligned}
\]

The factorial coefficients \(A_n=(2n+1)!u_n\) and \(B_n=(2n+1)!v_n\)
satisfy

\[
\begin{aligned}
A_n&=\sum_{i=0}^{n-1}\binom{2n}{2i+1}
 q^{2i+1}r^{2(n-1-i)+1}A_iB_{n-1-i},\\
B_n&=\sum_{i=0}^{n-1}\binom{2n}{2i+1}
 a^{2i+1}b^{2(n-1-i)+1}A_iB_{n-1-i}.
\end{aligned}
\]

Thus \(A_n,B_n\) are polynomials with nonnegative integer coefficients
in the four scales. These formulas are exact algebraic constructions.
Finite checks of them do not replace the analytic proof.

\hypertarget{v11-coupled_growth_full-the-characteristic-exponent-and-its-positive-vector}{%
\subsection{The characteristic exponent and its positive
vector}\label{v11-coupled_growth_full-the-characteristic-exponent-and-its-positive-vector}}

For \(s>0\), let

\[
P(s)=\begin{pmatrix}q^s&r^s\\a^s&b^s\end{pmatrix}.
\]

There is a unique \(\rho>0\) with

\[
\boxed{(1-q^\rho)(1-b^\rho)=(ar)^\rho.}
\]

Indeed, the left side strictly increases continuously from zero to one,
while the right side strictly decreases continuously from one to zero.
This equation is precisely \(\det(I-P(\rho))=0\). Define

\[
\widetilde w=(r^\rho,1-q^\rho)^T,
\qquad
w=\frac{\widetilde w}{\max(r^\rho,1-q^\rho)}.
\]

Both coordinates of \(w\) are positive, \(\max(w_1,w_2)=1\), and direct
multiplication gives

\[
P(\rho)w=w.
\]

The other eigenvalue is \(q^\rho+b^\rho-1\), whose absolute value is
strictly less than one. Consequently

\[
\boxed{\operatorname{spr}P(\rho)=1.}
\]

Conversely, a positive matrix of this form having spectral radius one
has eigenvalue one and satisfies the determinant equation. This proves
the proposed spectral characterization without requiring a numerical
eigensolver. A positive left eigenvector is

\[
\ell=(a^\rho,1-q^\rho)^T,\qquad \ell^TP(\rho)=\ell^T.
\]

The notation \(w\) for the growth vector is deliberately different from
\(v_n\), the coefficients of \(V\).

\hypertarget{v11-coupled_growth_full-exact-order-and-explicit-positive-type-bounds}{%
\subsection{Exact order and explicit positive type
bounds}\label{v11-coupled_growth_full-exact-order-and-explicit-positive-type-bounds}}

For \(R>0\), positivity gives

\[
M_U(R)=\max_{|z|=R}|U(z)|=U(R),\qquad
M_V(R)=V(R).
\]

Set \(F_1=U\), \(F_2=V\), and \(H_i(R)=\log F_i(R)\). Since
\(F_i(R)\ge R\), these logarithms are positive for \(R>1\). The type
convention used here is

\[
\underline\tau_i=\liminf_{R\to\infty}\frac{\log M_{F_i}(R)}{R^\rho},
\qquad
\overline\tau_i=\limsup_{R\to\infty}\frac{\log M_{F_i}(R)}{R^\rho}.
\]

\textbf{Theorem 2.} Both functions have order \(\rho\), positive lower
type and finite upper type. More precisely, the explicit constants below
satisfy

\[
0<c_-w_i\le\underline\tau_i\le\overline\tau_i\le c_+w_i<\infty,
\qquad i=1,2.
\]

Moreover,

\[
\lim_{R\to\infty}\frac{\log\log M_{F_i}(R)}{\log R}=\rho.
\]

This radial exponent limit is stronger than the limsup used to define
order. It does not assert that either type has a limit or that growth is
uniform in the angular variable.

\hypertarget{v11-coupled_growth_full-explicit-lower-comparison}{%
\subsubsection{Explicit lower
comparison}\label{v11-coupled_growth_full-explicit-lower-comparison}}

Let

\[
m=\min(q,r,a,b),\qquad Q=\max(q,r,a,b),
\]

\[
C=\frac{Q}{1-Q},\qquad B=1+\frac2m,
\qquad c_-=\frac{\log2}{(B+C)^\rho}.
\]

We prove the stronger statement

\[
\boxed{H_i(R)\ge c_-w_i(R+C)^\rho\quad(R\ge2,\ i=1,2).}
\]

On \(2\le R\le B\), its right side is at most \(\log2\), because
\(w_i\le1\), whereas \(H_i(R)\ge\log R\ge\log2\).

For the first equation, integrate only over the final unit interval to
obtain, when \(R\ge B\),

\[
U(R)\ge\int_{R-1}^R U(qt)V(rt)\,dt
\ge U(q(R-1))V(r(R-1)).
\]

The second equation gives the same inequality with \((q,r)\) replaced by
\((a,b)\) and \(U(R)\) replaced by \(V(R)\). All four child radii are at
least two. Every scale \(p\in\{q,r,a,b\}\) satisfies

\[
p(R-1)+C\ge p(R+C),
\]

because \((1-p)C\ge p\). If the comparison already holds at the child
radii, then

\[
\begin{aligned}
H_1(R)&\ge c_-\{w_1[q(R-1)+C]^\rho
                     +w_2[r(R-1)+C]^\rho\}\\
&\ge c_-(q^\rho w_1+r^\rho w_2)(R+C)^\rho
=c_-w_1(R+C)^\rho.
\end{aligned}
\]

The second row of \(P(\rho)w=w\) proves the second component in exactly
the same way.

For a formal propagation, put

\[
d=(1-Q)B+Q>0.
\]

Every child of \(R\ge B\) is at most \(Q(R-1)\le R-d\). Simultaneous
induction over the intervals \([B+kd,B+(k+1)d]\), starting with the
established seed interval, therefore proves both comparisons for every
\(R\ge2\). Continuity handles shared endpoints. The argument uses only
monotonicity of positive powers and works unchanged when \(0<\rho<1\).

\hypertarget{v11-coupled_growth_full-explicit-upper-comparison}{%
\subsubsection{Explicit upper
comparison}\label{v11-coupled_growth_full-explicit-upper-comparison}}

Set

\[
\Delta=\log(1/m),\qquad
D=\max\left(\log\frac1{qr},\log\frac1{ab}\right)+\log2.
\]

Choose any integer \(N\ge0\) satisfying \(Q^{-N}\ge3/m\), and define

\[
W_N=(2^{N+1}-1)\log2
 +(2^{N+1}-N-2)\log(1/Q),
\]

\[
S=W_N+1+\Delta+D,\qquad
c_+=\frac{e^{-\rho}S}{\min(w_1,w_2)}.
\]

These constants require only the parameters and the uniquely defined
root \(\rho\). They require no evaluation of an unknown solution.

To bound a seed interval, the coefficient majorant from Theorem 1 gives
\(F_i(1)\le2\). Put \(R_k=Q^{-k}\), \(Z_0=2\), and

\[
Z_k=Q^{-k}(1+Z_{k-1}^2).
\]

Monotonicity of the positive functions and the integral equations give
\(F_i(R_k)\le Z_k\) simultaneously. As \(Z_k\ge1\),

\[
\log Z_k\le k\log(1/Q)+\log2+2\log Z_{k-1}.
\]

Summing this elementary recurrence gives \(\log Z_N\le W_N\). Therefore
\(H_i(e^t)\le W_N\) throughout \(1\le t\le1+\Delta\), since
\(e^t\le e/m<3/m\le R_N\).

Now let

\[
G_i(t)=H_i(e^t)+t+D.
\]

For \(t\ge1+\Delta\), all child radii are at least \(e\), so their
function values are at least one. The full integral gives

\[
H_1(e^t)\le t+\log2
 +H_1(e^{t-\log(1/q)})+H_2(e^{t-\log(1/r)}).
\]

The definition of \(D\) converts this into

\[
G_1(t)\le G_1(t-\log(1/q))+G_2(t-\log(1/r)).
\]

Similarly,

\[
G_2(t)\le G_1(t-\log(1/a))+G_2(t-\log(1/b)).
\]

On the seed interval, \(G_i(t)\le S\le c_+w_i e^{\rho t}\). The Perron
equation propagates this comparison: substituting the bounds at the
children in the first inequality gives
\(c_+e^{\rho t}(q^\rho w_1+r^\rho w_2)=c_+w_1e^{\rho t}\), and the
second row works likewise. All logarithmic delays belong to
\([\log(1/Q),\Delta]\). Induction in intervals of length
\(\log(1/Q)>0\), after the seed interval, therefore proves

\[
\boxed{H_i(R)\le c_+w_iR^\rho-\log R-D\quad(R\ge e).}
\]

The lower and upper comparisons imply the asserted positive, finite
types. Taking logarithms of these bounds for sufficiently large \(R\)
gives the claimed radial exponent limits and hence exact orders.
\(\square\)

\hypertarget{v11-coupled_growth_full-the-two-type-vectors-have-exactly-the-perron-proportions}{%
\subsection{The two type vectors have exactly the Perron
proportions}\label{v11-coupled_growth_full-the-two-type-vectors-have-exactly-the-perron-proportions}}

The explicit constants above are conservative. There is nevertheless an
exact statement about the relative types of the two components.

\textbf{Theorem 3.} There exist positive finite constants
\(\underline\tau\le\overline\tau\) such that

\[
\boxed{\begin{aligned}
(\underline\tau_1,\underline\tau_2)^T&=\underline\tau\,w,\\
(\overline\tau_1,\overline\tau_2)^T&=\overline\tau\,w.
\end{aligned}}
\]

In particular,

\[
\frac{\underline\tau_1}{\underline\tau_2}
=\frac{\overline\tau_1}{\overline\tau_2}
=\frac{r^\rho}{1-q^\rho}
=\frac{1-b^\rho}{a^\rho}.
\]

Either both ordinary type limits exist or neither does. When they exist,
their vector is a positive multiple of \(w\).

\textbf{Proof.} Use the final-unit integral inequality from the lower
proof and divide by \(R^\rho\). Taking the liminf gives

\[
\underline t\ge P(\rho)\underline t,
\qquad
\underline t=(\underline\tau_1,\underline\tau_2)^T.
\]

Here \(p(R-1)\) tends to infinity continuously and
\([p(R-1)]^\rho/R^\rho\to p^\rho\), which justifies each scaled liminf.
The liminf of a sum is at least the sum of the liminfs. All quantities
are positive and finite by Theorem 2.

The full-integral inequality

\[
H_1(R)\le\log R+\log2+H_1(qR)+H_2(rR)
\]

and its second-row counterpart imply, on taking the limsup,

\[
\overline t\le P(\rho)\overline t.
\]

The terms \((\log R+\log2)/R^\rho\) vanish. Dotting either componentwise
inequality with the strictly positive left eigenvector \(\ell\) from
Section~\ref{v11-coupled_growth_full-the-characteristic-exponent-and-its-positive-vector} shows that its nonnegative difference vector has zero positive
weighted sum. Every coordinate of that difference is therefore zero.
Thus both type vectors are fixed by \(P(\rho)\). Its eigenvalue-one
space is one dimensional, so each is a positive multiple of \(w\).

If one component has an ordinary type limit, its lower and upper types
are equal. Because that coordinate of \(w\) is positive, this forces
\(\underline\tau=\overline\tau\), and the other component also has a
type limit. \(\square\)

This theorem does not assert a limit for \(\log U(R)/\log V(R)\) when
type oscillations persist. It also does not identify the scalar values
of either type or prove their equality.

\hypertarget{v11-coupled_growth_full-reductions-a-worked-unequal-pair-and-parameter-structure}{%
\subsection{Reductions, a worked unequal pair, and parameter
structure}\label{v11-coupled_growth_full-reductions-a-worked-unequal-pair-and-parameter-structure}}

\textbf{Exact scalar reduction.} The components coincide if and only if
the unordered pairs of scales coincide:

\[
U=V\quad\Longleftrightarrow\quad\{q,r\}=\{a,b\}.
\]

For sufficiency, the scalar solution \(T_{q,r}\) from
\path{routes/asymmetric_dilation.md} makes \((T_{q,r},T_{q,r})\)
solve the coupled equations. Theorem 1 gives uniqueness. For necessity,
equality at degree three gives \(qr=ab=p>0\). Equality at degree five
then gives \(q^2+r^2=a^2+b^2\). The two positive unordered pairs have
the same product and the same sum \(\sqrt{q^2+r^2+2p}\), so they
coincide. This characterizes equality of the normalized components. A
stronger assertion about all possible changes of variables requires its
own nonreduction theorem.

On either scalar reduction locus, the characteristic equation is
equivalent to \(q^\rho+r^\rho=1\), recovering the scalar order theorem.
At \(q=r=a=b=1/2\),

\[
U(z)=V(z)=\sqrt2\sinh(z/\sqrt2),\qquad \rho=1.
\]

\textbf{A genuinely unequal order-one pair.} Set

\[
(q,r,a,b)=\left(\frac15,\frac25,\frac35,\frac7{10}\right).
\]

The determinant identity at \(\rho=1\) is
\((1-1/5)(1-7/10)=(3/5)(2/5)=6/25\). Thus the exact order is one, and

\[
w=(1/2,1)^T.
\]

The first terms are

\[
\begin{aligned}
U(z)&=z+\frac2{75}z^3+\frac{88}{234375}z^5
 +\frac{18814}{8544921875}z^7+O(z^9),\\
V(z)&=z+\frac7{50}z^3+\frac{8211}{1250000}z^5
 +\frac{85664601}{625000000000}z^7+O(z^9).
\end{aligned}
\]

Both lower and upper exponential types of \(V\) are exactly twice the
corresponding types of \(U\). The row sums of \(P(1)\) are \(3/5\) and
\(13/10\). Consequently neither an unweighted average row condition nor
the maximum of the two scalar row orders yields this example's coupled
order. The positive eigenvector accounts for the exchange of growth
between the components.

For this example the elementary lower constants are \(m=1/5\),
\(Q=7/10\), \(C=7/3\), \(B=11\), and \(c_-=3\log2/40\). Hence, for every
\(R\ge2\),

\[
\log U(R)\ge\frac{3\log2}{80}\left(R+\frac73\right),\qquad
\log V(R)\ge\frac{3\log2}{40}\left(R+\frac73\right).
\]

\textbf{Dependence of the order on the scales.} The determinant equation
shows that the order depends on the self scales \(q,b\) and the product
\(ar\) of the two cross scales. Changing \(a\) and \(r\) while keeping
their product fixed leaves the order fixed, provided both remain in
\((0,1)\). The Perron vector and the normalized functions generally
change.

For any prescribed \(s>0\) and any \(q,r\in(0,1)\), choices satisfying

\[
a^s r^s=(1-q^s)(1-b^s),\qquad 0<a,b<1,
\]

lie on the order-\(s\) surface. For instance, one may select \(q,r,b\)
with \((1-q^s)(1-b^s)<r^s\), then define

\[
a=\frac{[(1-q^s)(1-b^s)]^{1/s}}{r}.
\]

This constructs exact order surfaces without fitting growth from
truncated series.

\hypertarget{v11-coupled_growth_full-proof-provenance-and-remaining-questions}{%
\subsection{Proof provenance and remaining
questions}\label{v11-coupled_growth_full-proof-provenance-and-remaining-questions}}

The decisive extension is from a scalar positive comparison to a
comparison weighted by a positive eigenvector. The fixed-width integral
controls the lower type at the exact exponent. A fixed proportional
interval \([\alpha R,R]\), with \(\alpha<1\), would only give a smaller
exponent through \(\alpha^sP(s)\). Letting \(\alpha\) tend to one
recovers exact order but does not, by itself, yield positive lower type.
This is a limitation of that alternative route, not a failure of the
theorem.

The complete proofs above close the target on the strictly positive
contraction domain. No computed example is used as a universal analytic
premise. The displayed rational Taylor coefficients were evaluated
directly from the exact recurrence with Python
\texttt{fractions.Fraction} as a transcription check. Their independent
formal residual checks belong to the release's exact computation record.

Remaining questions include the existence and value of ordinary type
limits, possible oscillations, zeros, angular growth, connection
formulas, and regimes with zero, negative, complex, or noncontracting
scales. Positivity and strict contraction are hypotheses of the results
here. This route makes no historical priority claim and no claim of a
Lean formalization.

Self-review coverage: this entire file. Routine checks: consistent
coefficient and eigenvector notation, domains for every logarithm and
power, seed and child interval endpoints, and both vector orientations.
No unresolved mathematical gap was identified in the author's
self-review. Independent review remains a separate record.
